\documentclass[11pt]{article}

\PassOptionsToPackage{numbers}{natbib}
\usepackage[final]{aaj2026}

\usepackage[T1]{fontenc}    % use 8-bit T1 fonts
\usepackage[spanish]{babel} % si tu instalación no trae babel-spanish, comenta esta línea
\usepackage{hyperref}       % hyperlinks
\usepackage{url}            % simple URL typesetting
\usepackage{booktabs}       % professional-quality tables
\usepackage{amsfonts}       % blackboard math symbols
\usepackage{nicefrac}       % compact symbols for 1/2, etc.
\usepackage{microtype}      % microtypography
\usepackage{xcolor}         % colors
\usepackage{verbatim}

\bibliographystyle{plainnat}

\title{Detección automática de series estadísticas\\
mediante comprensión de lenguaje natural:\\
aplicación al Banco Central de Chile}

\author{%
  Sebastián Mena Aliaga\\
  División de Estadísticas y Datos\\
  Banco Central de Chile \\
  \texttt{smena@bcentral.cl}\\
\And
 Hernán Fernández Carné\\
  División de Estadísticas y Datos\\
  Banco Central de Chile \\
  \texttt{hfernandez@bcentral.cl}\\
}

\begin{document}

\maketitle

\begin{abstract}
    El acceso a estadísticas económicas publicadas por bancos centrales requiere
    típicamente que el usuario conozca el catálogo de datos disponible, lo que
    representa una barrera para economistas y público general. En este trabajo se
    propone un sistema de comprensión de consultas en lenguaje natural que detecta
    de forma conjunta entidades e intenciones para identificar automáticamente la o
    las series estadísticas solicitadas y el modo en que deben ser entregados sus valores,
    aplicado a las series de la Base de Datos Estadísticos del Banco Central de Chile.
    Las entidades y clasificaciones detectadas son normalizadas y mapeadas mediante
    un catálogo estructurado hacia las series correspondientes. 
    \begin{comment}    
    El sistema es evaluado sobre
    consultas del PIB e IMACEC comparando tres enfoques: TF-IDF con SVM, mDeBERTa
    ajustado al dominio y FLAN-T5 como referencia generativa. Los resultados
    muestran que el enfoque propuesto alcanza una precisión de [X] en identificación
    de series y [X] en clasificación de intención, con potencial de extensión a
    otros indicadores y catálogos estadísticos institucionales.   

    == SE PODRÍA CAMBIAR CON LA COMPARACIÓN DE UN MODELO LLM CON MCP ==
    \end{comment}
\end{abstract}

\section{Introducción}

Los bancos centrales y organismos estadísticos publican vastos catálogos de series
de datos económicos que constituyen insumos fundamentales para el análisis
macroeconómico, la toma de decisiones de política pública y el seguimiento de la
actividad económica. Sin embargo, el acceso a estas series exige que el usuario
conozca con precisión la estructura interna del catálogo: qué indicadores existen,
cómo se denominan, qué desagregaciones están disponibles y bajo qué criterios se
organizan. Esta barrera limita el aprovechamiento de información que, en principio,
es pública y de alto valor.

El procesamiento de lenguaje natural (PLN) ofrece una vía para reducir esta
fricción, permitiendo que los usuarios expresen sus consultas de forma libre y que
un sistema sea capaz de interpretarlas y traducirlas al lenguaje del catálogo
institucional. No obstante, este problema va más allá de una simple búsqueda por
palabras clave: una consulta como \textit{``¿cuánto creció el comercio en la Región
de Valparaíso el último trimestre?''} contiene múltiples dimensiones de información
---el indicador, la actividad económica, la región, el período y el tipo de cálculo
requerido--- que deben ser identificadas y normalizadas de forma conjunta para
recuperar correctamente la serie solicitada.

\begin{comment}  

== MODIFICAR PARA QUE EL OBJETIVO SEA INCLUIR LA COMPARACIÓN DE MODELOS ==

En este trabajo se propone un sistema de comprensión de consultas en lenguaje
natural orientado a dos objetivos simultáneos: detectar \textit{qué serie
estadística} está siendo solicitada y determinar \textit{de qué forma} se requieren
sus datos. El sistema se basa en una arquitectura JointBERT multitarea que realiza
extracción de entidades y clasificación de intención de forma conjunta,
complementada con reglas de normalización y un diccionario estructurado para el
mapeo final hacia el catálogo del Banco Central de Chile (BCCh). El enfoque es
evaluado sobre series del Producto Interno Bruto (PIB) y el Índice Mensual de
Actividad Económica (IMACEC).

Las principales contribuciones de este trabajo son las siguientes:

\begin{itemize}
    \item Se formula el problema de recuperación de series estadísticas
    institucionales como una tarea conjunta de \textit{slot filling} y clasificación
    de intención sobre consultas en lenguaje natural.

    \item Se propone una arquitectura multitarea basada en JointBERT adaptada al
    dominio de estadísticas macroeconómicas, con entidades específicas al contexto
    del BCCh.

    \item Se construye un conjunto de datos anotado de consultas en lenguaje natural
    sobre series del PIB e IMACEC, junto con un esquema de normalización y mapeo
    hacia el catálogo oficial.

    \item Se evalúa el sistema y se muestra que el enfoque híbrido neural-simbólico
    es efectivo para este tipo de recuperación en dominios acotados.
\end{itemize}

\end{comment}

\section{Revisión del estado del arte}
\label{sec:estado-del-arte}

Investigadores del Banco Central Europeo (ECB) desarrollaron una interfaz de chatbot
utilizando técnicas de LLM y RAG. El procesamiento se basa en varios pasos
secuenciales. El primero, un procesamiento de la consulta que consiste en
estandarizarla usando un modelo LLM con un \textit{system prompt} preparado para
ello, con la finalidad de optimizarla para el \textit{data retrieval}. Luego, se
realiza una partición de la consulta en diferentes partes de la información. El
siguiente paso consiste en el \textit{data retrieval} para obtener las series
relevantes, usando \textit{embeddings} para almacenar las series y recuperarlas
mediante FAISS. Los investigadores solo incluyeron las 1.000 series más demandadas.
Las series obtenidas son evaluadas por un agente que juzga la relevancia y calidad
de la respuesta generada en contraste con la pregunta original; si es satisfactorio
el sistema procede a construir la respuesta específica, en caso contrario, se
realiza una segunda búsqueda que representa un subconjunto de aproximadamente
300.000 series. Si sigue insatisfactorio se concluye que no se puede responder a la
pregunta. Posteriormente, se obtiene la metadata de las series obtenidas y pasa por
el agente de generación de SQL para obtener los datos relevantes. Finalmente, un
agente redacta la respuesta al usuario.\footnote{\url{https://www.bis.org/ifc/publ/ifcb67_36.pdf}}

\section{La taxonomía de la BDE}
\label{sec:taxonomia_bde}

El presente trabajo se circunscribe a las estadísticas de cuentas nacionales del
Banco Central de Chile, específicamente al Producto Interno Bruto (PIB) y al
Indicador Mensual de Actividad Económica (Imacec), difundidas a través de la Base de
Datos Estadísticos (BDE). La BDE se organiza en dos niveles jerárquicos. El cuadro
constituye la unidad de publicación: un conjunto de series temporales que comparten
indicador, compilación de referencia o base metodológica, tratamiento estacional y
forma de presentación ---nivel, contribución o participación---; en el caso del PIB
comparten además la valoración, a precios corrientes o en volumen a precios del año
anterior encadenado, distinción que no aplica al Imacec por tratarse de un índice. La
serie constituye la unidad de dato: una sucesión temporal de observaciones
identificada por un código único, correspondiente a una fila del cuadro. Cada cuadro
presenta el indicador desagregado según una clasificación determinada, de modo que
cada una de sus series corresponde a una categoría de esa clasificación: en el Imacec
dicha clasificación suele ser la actividad económica, mientras que el PIB se
desagrega por actividad, región, gasto o ingreso. Cabe destacar que las
combinaciones de dimensiones efectivamente publicadas constituyen un subconjunto
disperso de las posibles: en el Imacec, por ejemplo, la contribución año/año se
publica sobre series originales y la contribución mes/mes sobre series
desestacionalizadas, sin que existan las combinaciones restantes.

El universo considerado comprende 280 cuadros y 4.497 series, de las cuales 1.503
corresponden a publicaciones vigentes y el resto a compilaciones de referencia y
bases metodológicas ya superadas. Se excluyeron los cuadros de deflactor del PIB,
agregados macroeconómicos, formación bruta de capital fijo, PIB e ingreso, y
versiones sucesivas de publicación. El Cuadro~\ref{tab:delimitacion} presenta la
composición del universo dentro del capítulo de Cuentas Nacionales.

\begin{table}[htbp]
\centering
\caption{Delimitación del universo según vigencia de publicación}
\label{tab:delimitacion}
\small
\begin{tabular}{l rrr rrr}
\toprule
& \multicolumn{3}{c}{Cuadros} & \multicolumn{3}{c}{Series} \\
\cmidrule(lr){2-4}\cmidrule(lr){5-7}
Familia & Vig. & Hist. & Total & Vig. & Hist. & Total \\
\midrule
\multicolumn{7}{l}{\textbf{Incluidos}} \\
\addlinespace[2pt]
\quad PIB regional por actividad económica  &  65 &  99 & 164 &  884 & 1.326 & 2.210 \\
\quad Gasto del PIB                         &  10 &  35 &  45 &  238 &   664 &   902 \\
\quad PIB por actividad económica           &   9 &  27 &  36 &  264 &   756 & 1.020 \\
\quad Imacec                                &   4 &  13 &  17 &   38 &    59 &    97 \\
\quad PIB por región                        &   5 &   9 &  14 &   76 &   160 &   236 \\
\quad PIB agregado                          &   1 &   3 &   4 &    3 &    10 &    13 \\
\addlinespace[2pt]
\quad \emph{Subtotal} & \textbf{94} & \textbf{186} & \textbf{280} & \textbf{1.503} & \textbf{2.975} & \textbf{4.478} \\
\addlinespace[5pt]
\multicolumn{7}{l}{\textbf{Excluidos}} \\
\addlinespace[2pt]
\quad Cuentas financieras por sector institucional & 50 &  80 & 130 & 1.417 & 2.215 & 3.632 \\
\quad Gasto regionalizado                          & 104 &  0 & 104 &   651 &     0 &   651 \\
\quad Stock de capital e inversión                 &  19 & 15 &  34 &   250 &   445 &   695 \\
\quad Operaciones no financieras por sector        &  11 & 23 &  34 &    89 &   219 &   308 \\
\quad Producto e ingreso nacional                  &   5 & 25 &  30 &    39 &   191 &   230 \\
\quad Deflactores                                  &   6 & 20 &  26 &   160 &   511 &   671 \\
\quad Ratios sectoriales                           &   5 & 10 &  15 &    31 &    53 &    84 \\
\quad Bases de datos de revisiones                 &  11 &  0 &  11 &   858 &     0 &   858 \\
\quad Imacec sectorial                             &   2 &  4 &   6 &    56 &   120 &   176 \\
\addlinespace[2pt]
\quad \emph{Subtotal} & \textbf{213} & \textbf{177} & \textbf{390} & \textbf{3.551} & \textbf{3.754} & \textbf{7.305} \\
\midrule
Total capítulo Cuentas Nacionales & 307 & 363 & 670 & 5.054 & 6.729 & 11.783 \\
\bottomrule
\end{tabular}
\end{table}

\subsection{Delimitación del ámbito de la consulta}

El primer objetivo es identificar si una pregunta pertenece al ámbito económico, si
se relaciona con los indicadores del PIB o el IMACEC publicados por el Banco Central
de Chile, y si la pregunta puede contestarse únicamente a partir de la información
oficial publicada sin necesidad de realizar análisis o interpretación de esta, con la
finalidad de resguardar la reputación del Banco y responder solo aquellas preguntas
que le corresponde contestar.

La literatura aborda este problema como detección de consultas fuera de dominio
(\textit{out-of-scope}) y como \textit{answerability}, y modela el rechazo como
incapacidad: el sistema no puede responder porque el dato no está, el esquema no lo
soporta o la intención no fue entrenada. En el caso aquí considerado, en cambio, el
rechazo obedece a un mandato institucional y no a una limitación de capacidad.

Las tres condiciones se desagregan en: ¿es del ámbito económico?, ¿se relaciona con
el PIB o el IMACEC?, y ¿se responde directamente a partir de lo publicado, sin
análisis? En consecuencia, se propone un clasificador de tres cabezas binarias que
etiqueten la pregunta según cada una de estas condiciones.

\section{Arquitectura del PIBERT}
\label{sec:pibert}

El Banco publica periódicamente las series del PIB y el Imacec organizadas por
cuadros. El segundo objetivo del sistema es identificar cuáles de estas series y
cuadros se mencionan, directa o indirectamente, en la pregunta. Un cuadro agrupa
un conjunto de series; una misma serie puede figurar en más de un cuadro cuando
el catálogo ofrece cortes alternativos del mismo dato, de modo que la resolución
devuelve un conjunto de pares cuadro--serie antes que una serie aislada.

La tarea de identificar correctamente los cuadros y series que subyacen a una
pregunta implica la capacidad de interpretarla a partir de un conocimiento
acabado del dominio. Se propone un proceso sistematizado de interpretación e
identificación utilizando herramientas de Procesamiento de Lenguaje Natural,
basado en la clasificación de etiquetas y el Reconocimiento de Entidades
Nombradas (NER) mediante un modelo de tipo JointBERT. La selección de las cabezas
clasificadoras y los tipos de entidad que participan en la identificación
conllevó la iteración, basada en prueba y error, de diferentes arquitecturas
capaces de generalizar y abarcar la mayor cantidad de casos posibles. La
arquitectura final se compone de ocho dimensiones enunciables en la consulta,
tres cabezas clasificadoras binarias, un catálogo estructurado de resolución y un
conjunto de valores por defecto.

\subsection{Tipos de entidad y normalización}
\label{sec:entidades}

Cada serie del catálogo queda descrita por una tupla de catorce dimensiones,
derivada de su identificador oficial. Sólo ocho de ellas admiten mención en una
consulta; son las que el extractor debe resolver, mediante reconocimiento de
entidades y normalización del span contra un vocabulario controlado. Cada tipo
de entidad tiene asociado un conjunto de valores canónicos, y la normalización
es la función que mapea las múltiples formas con que una mención puede aparecer
en el texto a un único valor de ese conjunto. Dos dimensiones adicionales no
identifican la serie sino la solicitud sobre ella:

\begin{verbatim}
Dimensiones de identidad (enunciables en la consulta)
INDICADOR       imacec, pib
FRECUENCIA      mensual, trimestral, anual
ESTACIONALIDAD  original, desestacionalizada
PRESENTACION    nivel, contribucion_aa, contribucion_pa, participacion
                [la ausencia de mencion se resuelve a nivel]
VALORACION      corrientes, encadenado, constantes
ACTIVIDAD       mineria, industria_manufacturera, comercio, ...   (45)
REGION          metropolitana, valparaiso, los_lagos, ...          (16)
GASTO           demanda_interna, consumo_total, fbcf, ...          (35)

Dimensiones de la solicitud (no identifican la serie)
PERIODO         fechas y expresiones temporales relativas
CALCULO         nivel, var_periodo_anterior, var_anual
\end{verbatim}

\medskip

\noindent Las seis dimensiones restantes de la tupla no dan lugar a mención en
el uso corriente y se resuelven fuera del modelo: la compilación de referencia
(2018, 2013, 2008, \dots) sólo se enuncia excepcionalmente y toma por defecto la
vigente; la vigencia de publicación, el empalme y la condición de información
histórica se fijan por valores por defecto (Sección~\ref{sec:defaults}); y la
apertura del cuadro y la región fija se derivan del propio catálogo. Los valores
marginales del catálogo ---tasas, series de tendencia-ciclo y la fila
extrarregional--- se conservan en la tabla de resolución pero se excluyen del
vocabulario del modelo.

Los vocabularios de ACTIVIDAD, REGION y GASTO no se exponen al modelo como
clases: el extractor etiqueta un span genérico del tipo correspondiente y la
normalización lo resuelve contra las categorías de los cuadros candidatos, de
modo que la granularidad reside en el catálogo y no en el clasificador. PERIODO
y CALCULO especifican el recorte temporal y la operación solicitada, y operan
como fuente indirecta de desambiguación sobre las dimensiones no enunciadas.

A modo de ilustración, en la consulta «cuánto creció el producto interno bruto
sin estacionalidad» se reconocen tres menciones ---«creció», «producto interno
bruto» y «sin estacionalidad»--- correspondientes a los tipos CALCULO, INDICADOR
y ESTACIONALIDAD, que la normalización resuelve respectivamente a los valores
canónicos \texttt{var\_anual}, \texttt{pib} y \texttt{desestacionalizada}. La
función es de muchos a uno y no reversible: las formas «sin estacionalidad»,
«desestacionalizado» y «serie ajustada» convergen en un mismo valor canónico,
pérdida de información deliberada que otorga al sistema invarianza frente al
registro en que se formula la consulta.

%%%% REVISAR !!!! %%%%
Tres de los vocabularios son condicionales. El de ACTIVIDAD depende del cuadro, dado que el Imacec emplea una clasificación agregada por actividad económica  mientras que los cuadros de PIB emplean desagregaciones considerablemente más finas. El de PRECIO depende de la compilación de referencia, puesto que la valoración real se expresa en volumen encadenado desde la referencia 2008 y en precios constantes de año base en las compilaciones anteriores. El de REGION también depende de la referencia, en tanto la división político-administrativa vigente varía entre compilaciones.

\subsection{Cabezas clasificadoras}
\label{sec:cabezas}

Las cabezas clasificadoras cumplen una función complementaria a la del reconocimiento de entidades: los spans capturan las menciones que refieren a una categoría y admiten normalización contra un vocabulario controlado, mientras que las cabezas resuelven las dimensiones que no dan lugar a una mención normalizable.

Esta complementariedad se manifiesta con particular claridad en el eje de desagregación. Los cuadros del Imacec y del PIB, seleccionados, se aperturan por actividad económica, región, actividad-región, o componentes del gasto, de modo que el span normalizado sobre dicho eje ---ACTIVIDAD, REGIÓN o GASTO--- constituye la pieza final de la resolución, una vez determinadas las restantes dimensiones. Cabe precisar que el nivel en que opera esa pieza depende del cuadro: en los cuadros desagregados por actividad, la mención selecciona una serie especifica; en la desagregación regional, en cambio, cada región puede ser una serie o constituir un cuadro independiente desagregado por actividad económica. Cuando la consulta no nombra categoría alguna pero la cabeza del eje se activa, como en «cuánto crecieron las actividades del Imacec el mes pasado», la resolución no apunta a una serie sino al cuadro en su totalidad.

Las cabezas se denotan \texttt{activ\_cls}, \texttt{reg\_cls} y \texttt{gast\_cls}, con valores en $\{0,1\}$, e indican si la consulta invoca la clasificación correspondiente en su conjunto. Las tres admiten activación simultánea, lo que permite representar los cuadros de doble desagregación por región y actividad, que constituyen la mayoría del catálogo considerado.

El eje de desagregación queda determinado por la unión de ambas señales ---la presencia de un span de ese tipo o la activación de la cabeza correspondiente---, mientras que su combinación fija la cardinalidad de la respuesta, según muestra el Cuadro~\ref{tab:cardinalidad}.

\begin{table}[htbp]
\centering
\caption{Cardinalidad de la respuesta según las señales sobre el eje de desagregación}
\label{tab:cardinalidad}
\begin{tabular}{ccl}
\toprule
Span & Cabeza & Resolución \\
\midrule
--    & 0 & fila agregada del cuadro \\
sí    & 0 & las categorías mencionadas \\
--    & 1 & el cuadro en su totalidad \\
sí    & 1 & el cuadro en su totalidad, con las categorías mencionadas destacadas \\
\bottomrule
\end{tabular}
\end{table}

\subsection{Catálogo de resolución}
\label{sec:catalogo}

Sea $\mathcal{C}$ el catálogo, donde cada cuadro $c \in \mathcal{C}$ se representa como una tupla completamente especificada $c = (c_1,\dots,c_k)$ sobre los vocabularios $V_1,\dots,V_k$ de las dimensiones descritas en la Sección~\ref{sec:entidades}. La tabla se deriva de la metadata de publicación sin construcción manual: el identificador de cada serie está compuesto por segmentos que codifican las dimensiones del cuadro al que pertenece, de modo que $\mathcal{C}$ se obtiene decodificando los identificadores del catálogo oficial. Esta propiedad hace el procedimiento transferible a cualquier oficina estadística que difunda sus series mediante claves estructuradas.

La consulta, en cambio, produce una tupla \emph{parcial}
$q \in \prod_i (V_i \cup \{\bot\})$, donde $\bot$ denota una dimensión que el
extractor no resolvió. Las posiciones no especificadas no restringen, de manera
que el conjunto candidato queda definido por
\[
  \mathcal{C}(q) \;=\; \{\, c \in \mathcal{C} \;:\;
  \forall i,\; q_i = \bot \;\lor\; q_i = c_i \,\}.
\]

La resolución concluye cuando $|\mathcal{C}(q)| = 1$. Si $|\mathcal{C}(q)| > 1$,
la dimensión de mayor cardinalidad en el conjunto residual señala el eje sobre el
cual formular una pregunta aclaratoria. Si $\mathcal{C}(q) = \emptyset$, la
combinación solicitada no existe en el catálogo, lo que constituye evidencia de
un error de extracción o de una consulta ajena al universo publicado.

La eficacia de este filtrado descansa en la dispersión del catálogo, que es
extrema: las once dimensiones de cuadro admiten $5{,}16 \times 10^{6}$
combinaciones, de las cuales el catálogo publica 275 tuplas distintas para 276
cuadros. Cabe observar que la proyección no es completamente inyectiva: un par
de cuadros comparte tupla ---el PIB trimestral y el PIB anual por región de la
referencia 2018--- por una inconsistencia de frecuencia en la metadata de
origen, que codifica ambos como trimestrales.
% TODO: cuadrar 276 con el Cuadro 1 (280).

\subsection{Resolución de dimensiones no enunciadas}
\label{sec:defaults}

Sobre el conjunto candidato operan dos mecanismos de naturaleza distinta.

El primero es la \emph{propagación de restricciones}, que colapsa aquellas
posiciones $\bot$ que toman un valor único a lo largo de $\mathcal{C}(q)$: la
dimensión queda determinada por la estructura del catálogo, sin que la consulta
la enuncie ni el modelo la infiera. Así, en «cuánto creció la economía el último
trimestre», el reconocimiento de la frecuencia trimestral excluye al Imacec
---que solo se publica con periodicidad mensual--- y fuerza la resolución hacia
el Producto Interno Bruto; de manera análoga, en «cuánto aceleró la economía», la
mención del cálculo respecto del período inmediatamente anterior determina que la
serie solicitada sea la desestacionalizada, dado que la variación intraperiodo
únicamente se publica sobre series ajustadas por estacionalidad. Estas
inferencias no requieren reglas explícitas ni entrenamiento adicional: emergen
del filtrado contra el conjunto de combinaciones efectivamente publicadas, cuya
dispersión codifica las convenciones metodológicas del compilador.

El segundo mecanismo es la asignación de \emph{valores por defecto} sobre las
posiciones que permanecen indeterminadas tras dicho colapso. Los defaults son
de tres clases, según cómo se imponen sobre el conjunto candidato
(Cuadro~\ref{tab:defaults}).

Las \emph{constantes} actúan como filtro previo a toda inferencia lingüística:
el sistema trabaja sobre las publicaciones vigentes de la compilación de
referencia actual, y sólo abre el universo histórico cuando la consulta nombra
una compilación distinta. Este filtro reduce el espacio de resolución de 276 a
91 cuadros.\footnote{Cifras de la exportación de la BDE de agosto de 2026; en la
publicación vigente todos los cuadros corresponden a la referencia 2018 y
ninguno a información histórica, de modo que la vigencia basta por sí sola para
la reducción.}
% TODO: cuadrar con el Cuadro 1 (280 / 94).

Las \emph{preferencias} son defaults blandos: se imponen únicamente cuando el
conjunto candidato contiene el valor preferido y no restringen en caso
contrario, de modo que nunca lo vacían. Su naturaleza responde a que la
convención no es uniforme en el catálogo. El empalme es el caso más claro: de
los 91 cuadros vigentes sólo 20 publican series empalmadas, y las
contribuciones y participaciones no lo hacen, por lo que imponerlo como
constante anularía la mayoría de las resoluciones; como preferencia, en cambio,
actúa sólo como desempate entre cuadros por lo demás equivalentes. Lo mismo
ocurre con la valoración, donde un mismo cuadro puede ofrecer la serie a
precios corrientes y en volumen encadenado, y la lectura convencional es la
segunda. Las preferencias se aplican en un orden fijo y cada una que se impone
queda registrada como supuesto, que se explicita en la respuesta entregada al
usuario.

La \emph{tabla condicionada}, por último, recoge convenciones de lectura cuyo
valor depende de otras dimensiones de la consulta: el tratamiento estacional se
resuelve según el indicador y el cálculo solicitado. A diferencia de las
anteriores, no está determinada por la estructura del catálogo: debe ser
mantenida institucionalmente, versionada y auditable. Cuando ambas alternativas
se encuentran publicadas y ninguna convención las discrimina, la tabla no
resuelve y la consulta se deriva a una pregunta aclaratoria sobre esa dimensión.

\begin{table}[htbp]
\centering
\caption{Valores por defecto sobre las dimensiones no enunciadas}
\label{tab:defaults}
\begin{tabular}{lll}
\toprule
Dimensión & Valor por defecto & Clase \\
\midrule
Vigencia               & vigente                        & constante \\
Información histórica  & no                             & constante \\
Referencia             & compilación 2018               & constante, salvo mención \\
Presentación           & nivel                          & preferencia \\
Valoración             & volumen encadenado             & preferencia \\
Empalme                & serie empalmada                & preferencia (desempate) \\
Estacionalidad         & según indicador y cálculo      & tabla condicionada \\
\bottomrule
\end{tabular}
\end{table}

La aplicación conjunta de estos mecanismos permite resolver consultas que no
enuncian dimensión alguna. La ausencia total de menciones, como en «cuánto creció
la economía», se resuelve hacia el Imacec en su serie original de frecuencia
mensual, compilación de referencia vigente y variación respecto de igual período
del año anterior, por tratarse de la lectura convencional de la pregunta en el
uso institucional.

\subsection{Catálogo de reglas de identificación}

Se define un catálogo que describe los cuadros y series como una tupla de los
valores posibles de las cabezas clasificadoras y las entidades normalizadas, desde
los casos más explícitos y evidentes hasta los casos implícitos que requieren un
conocimiento más especializado de interpretación.

\begin{verbatim}
# Identificación del indicador
CONJ_IMACEC: (INDIC=[] y FREQ=([] o ["m"])) o
             (INDIC=["imacec"] y FREQ=(["m"] o ["t"] o ["a"]))
CONJ_PIB:    (INDIC=([] o ["pib"]) y FREQ=(["t"] o ["a"]))

# Identificación de la estacionalidad
CONJ_EST:  (EST=[] y CALC=([] o ["pct"])) o
           (EST=["est"] y CALC=([] o ["pct"] o ["ypct"]))
CONJ_DEST: (EST=[] y CALC=([] o ["pct"])) o
           (EST=["dest"] y CALC=([] o ["pct"] o ["ypct"]))
\end{verbatim}

\noindent El tipo de cálculo informa directamente sobre el tipo de estacionalidad
solicitada. Un cuadro se identifica con la combinación de las entidades y
clasificadores INDIC, FREQ, DESC, EST, PREC, activ\_cls, reg\_cls, gast\_cls e
hist\_cls. A continuación se presentan cuatro ejemplos de reglas:

\begin{verbatim}
Imacec, volumen a precios del año anterior encadenado,
series empalmadas (promedio 2018=100)
  NER:     CONJ_IMACEC y CONJ_EST y DESC=[] y PREC=([] o ["enc"])
  Clasif.: activ_cls=1 ; reg_cls=0 ; gast_cls=0

Imacec, contribución porcentual respecto de igual periodo
del año anterior, referencia 2018
  NER:     CONJ_IMACEC y DESC=["contr"] y ESTAC=([] o ["est"])
           y PREC=([] o ["enc"])
  Clasif.: activ_cls=1 ; reg_cls=0 ; gast_cls=0

Producto interno bruto, referencia 2018
  NER:     CONJ_PIB y DESC=[] y ESTAC=([] o ["est"] o ["dest"])
           y PREC=([] o ["enc"] o ["corr"])
  Clasif.: activ_cls=0 ; reg_cls=0 ; gast_cls=0

Participación por clase de actividad económica,
series empalmadas, referencia 2018
  NER:     CONJ_PIB y DESC=["part"] y ESTAC=([] o ["est"])
           y PREC=([] o ["enc"] o ["corr"])
  Clasif.: activ_cls=0 ; reg_cls=0 ; gast_cls=0
\end{verbatim}

\begin{ack}
Use unnumbered first level headings for the acknowledgments. All acknowledgments
go at the end of the paper before the list of references. Moreover, you are required
to declare funding (financial activities supporting the submitted work) and
competing interests (related financial activities outside the submitted work).
\end{ack}

\medskip

\bibliography{biblio}

\appendix

\section{Appendix / supplemental material}

Optionally include supplemental material (complete proofs, additional experiments
and plots) in appendix.

\end{document}